% Supplement to the existing scaled-cluster result, not a new standalone paper.
\section{A finite rectangular partition of disjoint lattice cubes}
\label{sec:explicit-cube-partition}
The linear partition for supplied disjoint cubes can be derived directly,
without importing a general fat-object or low-density theorem. We use the
classical compressed cube-splitting strategy of de Berg, UU-CS-1998-12,
Section~2.1 and Lemma~2.1, and prove the required special-case geometry and
all discrete/energy adapters below. This is not an original general BSP
method. No time-complexity, measured implementation or Lean claim is made.

Fix $D\geq3$ and an integer $n\geq1$. In the common root
$X=\{0,\ldots,n-1\}^D$, take pairwise disjoint nonempty integer cubes
\[
 S_k=\prod_i\{b_{ki},\ldots,b_{ki}+L_k-1\},\qquad 1\leq k\leq J.
\]
Set
\[
 Q_k=\prod_i[b_{ki}-1/4,b_{ki}+L_k-3/4].
\]
These are positive-side closed solid cubes, including $L_k=1$. They are
strictly disjoint: disjoint integer intervals in a separating coordinate
leave a real gap at least $1/2$. Their integer points are exactly $S_k$,
and each such point is interior to $Q_k$. All corners of the $Q_k$ lie
strictly inside $[-1,n]^D$.

\subsection{Proper-intersection corner bound}
Let $C$ be any axis-aligned cube of side $s>0$ with no corner of any $Q_k$
in its interior. Then at most $2^D$ of the $Q_k$ meet $\operatorname{int}C$.
Indeed, if $Q_k$ of side $a<s$ properly meets $C$, each of its coordinate
intervals has an endpoint in the interior of the corresponding interval
of $C$. Choosing these endpoints gives a corner of $Q_k$ in
$\operatorname{int}C$, a contradiction. Thus every properly meeting $Q_k$
has $a\geq s$. In each coordinate its interval then contains at least
one endpoint of the interval of $C$; otherwise it would lie strictly
inside that shorter interval. Hence $Q_k$ contains a corner of $C$.
Strict disjointness makes these corner choices injective, proving the bound.

The qualifier ``properly'' is essential. Closed boundary-only contact
must not be counted here. For $C=[0,1]^D$ and any integer $K\geq1$, put
$\epsilon=1/(4K+4)$ and
\[
 Q_j=[-\epsilon,0]\times[(2j+1)\epsilon,(2j+2)\epsilon]
       \times[\epsilon,2\epsilon]^{D-2},\qquad0\leq j<K.
\]
These $K$ strictly disjoint small cubes touch $C$ but none meets its
interior or has a corner in its interior. Thus the corresponding
closed-contact bound would be false. This boundary guard does not assert
that de Berg's theorem is false; our construction below explicitly uses
open-cell intersections and interior lattice points.

\subsection{A multiple-object stress family}
The corner-free property does not imply that a query cube meets only one
input cube. This can fail even for the quarter-offset solids of equal-side
integer cubes, not merely for arbitrary real objects. For $D\geq2$, put
$n=4$, $C=[1,2]^D$, and, for $\sigma\in\{0,1\}^{D-1}$, set
\[
 b_\sigma=(2\sigma_1,\ldots,2\sigma_{D-1},1),\quad L_\sigma=2,
 \qquad S_\sigma=\prod_{i=1}^D\{b_{\sigma i},b_{\sigma i}+1\}.
\]
The $2^{D-1}$ integer cubes are pairwise disjoint: distinct binary words
have a coordinate with the disjoint integer intervals $\{0,1\}$ and
$\{2,3\}$. Their solids have, in that coordinate, intervals
$[-1/4,5/4]$ and $[7/4,13/4]$, with a gap $1/2$.
In the last coordinate every solid has interval $[3/4,9/4]$.
Consequently every solid properly meets $C$, while no solid corner is
interior to $C$: every such corner has last coordinate $3/4$ or $9/4$.
Thus a universal proper-occupancy bound for corner-free query cubes in
this exact lattice class must be at least $2^{D-1}$. This neither proves
sharpness of $2^D$ nor asserts that this particular $C$ occurs as a leaf
of our root-based recursion. It rules out replacing the occupancy
argument by a one-object assumption. The finite unit139 baseline leaves
all had occupancy at most one, so those runs did not test this phenomenon.
\subsection{Constructing linearly many vertex-free rectangles}
Let $V$ be the set of all corners of the $Q_k$, so $N=|V|=2^DJ$.
For $J=0$ use the single root rectangle and skip this construction.
For $J>0$ begin with the cube $[-1,n]^D$. At a current cube $C$ keep
exactly $V_C=V\cap\operatorname{int}C$.
If $V_C$ is empty, stop at that cube. Otherwise consider its $2^D$ open
half-side children.

If the points are not all in one open child, perform all midpoint cuts
and recurse on the nonempty child point sets. Points on the cuts are
removed from subsequent interior sets. If all points lie in one open
child, reflect and translate coordinates so the common corner is zero
and $C=[0,s]^D$. Put
\[
 a=\max\{x_i:x\in V_C,\ 1\leq i\leq D\},\qquad C'=[0,a]^D.
\]
Then $0<a<s/2$. Successively cut at the $D$ upper facet planes of $C'$,
producing $C'$ and the $D$ boxes
\[
 E_i=[0,a]^{i-1}\times[a,s]\times[0,s]^{D-i},\qquad1\leq i\leq D.
\]
Use disjoint boundary ownership for these boxes. Only $C'$ can contain
points of $V$ in its interior, and at least one point of $V_C$ lies on
its boundary. Stop at the other boxes and recurse inside $C'$ if needed.
Every intermediate or terminal cell is full-dimensional; a boundary
ownership convention never creates a separate hyperplane cell.

Here is a direct finite termination and size proof. Among active cells,
let $n_C=|V_C|>0$ and use the nonnegative integer potential
\[
 \Psi=\sum_{C\ {\rm active}}(2n_C-1).
\]
Initially $\Psi=2N-1$. If a processed cell loses $r$ points to new cut
boundaries and has $t$ nonempty child point sets, the decrease is
$2r+t-1$. In the shrinking step $r\geq1$. In a midpoint step either
$r\geq1$ or $t\geq2$. Thus every event decreases $\Psi$ by at least one,
so there are at most $2N-1$ events. Each event creates at most $2^D-1$
additional leaves, since a shrinking step creates only $D$ extra leaves.
The number $M$ of terminal rectangles satisfies
\begin{equation}\label{eq:explicit-first-size}
 M\leq1+(2N-1)(2^D-1)
 \leq2^{D+1}(2^D-1)J.
\end{equation}
Tracking exactly $V\cap\operatorname{int}C$ ensures that every terminal
rectangle has no corner of an input cube in its interior. This argument
is a construction and a decreasing integer potential, not a supplied
partition certificate. It does not prove a fast implementation bound.

\subsection{A uniform occupancy bound for the terminal rectangles}
Terminal cubes meet at most $2^D$ input cubes properly, by the corner bound.
A terminal noncube has the form $E_i$ above after a coordinate reflection.
It is contained in the slab $[0,s]^{i-1}\times[a,s]\times[0,s]^{D-i}$.
Cover this slab by $2^{D-1}$ closed cubes of side $s-a$: in coordinate $i$
use $[a,s]$, and in every other coordinate choose either $[0,s-a]$ or
$[a,s]$. These intervals cover $[0,s]$ since $a<s/2$.
Every covering cube lies inside the parent cube and outside the interior
of $C'$, and its interior contains no point of $V$. The corner bound
therefore applies separately to each covering cube.

If an input cube properly meets $E_i$, the intersection contains an open
set. A point of that set avoiding the finitely many covering-facet planes
belongs to the interior of a covering cube. Thus all proper intersections
are counted by the cover, even if some particular point lies on a cover
boundary. Consequently each terminal rectangle properly meets at most
\begin{equation}\label{eq:explicit-occupancy}
 T_D=2^{D-1}2^D=2^{2D-1}
\end{equation}
input cubes. The covering argument, including its boundary qualification,
is proved here; no arbitrary-box unclutteredness is assumed.

\subsection{From open geometry to an actual integer partition}
For each terminal rectangle list the at most $T_D$ cubes meeting its
interior. Refine by all their facet coordinates. There are at most $2T_D$
cut coordinates per axis, so at most $(2T_D+1)^D$ refined rectangles per
terminal cell. All cuts are axis-parallel. Now assign every splitting
hyperplane to one child consistently, intersect with $X$, and delete empty
integer sets. The result is a disjoint cover of $X$ by integer rectangles.
Taking integer points and deleting empty sets does not increase the count.
More explicitly, use the uniform convention $[l_i,h_i)$ in every axis
of every cell. A coordinate split at $t$ then assigns the plane to the
upper child. The two intervals $[l,t)$ and $[t,h)$ partition $[l,h)$
exactly, even when $t$ is an integer. By induction through every cut,
these half-open cells partition $[-1,n)^D$; intersecting with $X$ gives
\[
 R=\prod_i\bigl(\{\lceil l_i\rceil,\ldots,\lceil h_i\rceil-1\}
                 \cap\{0,\ldots,n-1\}\bigr).
\]
This formula is also valid for negative or integer endpoints. An empty
axis makes the entire integer cell empty; no lattice point is assigned
twice or lost when those cells are deleted. The outer upper face at $n$
contains no point of $X$.

To check membership rigorously, suppose an assigned integer point $x$
lies in $S_k$. It is interior to $Q_k$, while its assigned terminal
rectangle has $x$ in its closure. A small ball about $x$ lies in $Q_k$
and meets the interior of that full-dimensional terminal rectangle.
Therefore $Q_k$ was in its list. None of the new facet coordinates is an
integer. Hence each refined integer rectangle is entirely inside a
listed $S_k$ or disjoint from every listed $S_k$. An unlisted cluster
cannot contain one of its integer points, by the same closure argument.
Strict disjointness then implies that each integer rectangle is inside
one cluster or outside every cluster.

In particular, with the explicit dimension-only constant
\begin{equation}\label{eq:explicit-final-size}
 A_D=2^{D+1}(2^D-1)(2^{2D}+1)^D,
 \qquad |\mathcal P|\leq A_D\max\{1,J\},
\end{equation}
we have derived the required partition. The constant is deliberately
loose and no sharpness or algorithm running-time claim is made.
For $n=1$ and for singleton input cubes all constructions remain valid.

\subsection{The actual coefficient and fixed-penalty consequence}
Suppose, in addition, that for fixed $c>0$ and $\beta\geq1$ the actual
positive scalar starting coefficients obey
\[
 \begin{gathered}
 a(x)\geq cL_k(L_k-1)/4\quad(x\in S_k),\\
 a(x)\geq cn(n-1)/4\quad(x\notin\cup_k S_k),\\
 a(x)\leq\beta L_k^2\quad(x\in S_k).
 \end{gathered}
\]
These coefficient/cube assumptions are retained, not inferred from the
critical moment. On an interval of $l$ integer vertices, pairwise variance,
path Cauchy--Schwarz and at most $l^2/4$ pairs through one cut give
\[
 \sum|u-\bar u|^2\leq\frac{l(l-1)}4\sum_{\rm edges}|du|^2.
\]
For completeness, product variance tensorization follows from the commuting
coordinate averages $P_i u=l_i^{-1}\sum_{t=0}^{l_i-1}u(x_1,\ldots,t,\ldots,x_D)$: the telescoping orthogonal decomposition of
$I-\prod_iP_i$ gives a sum bounded by
$\sum_i\|u-P_i u\|_2^2$. Applying the interval bound on its fibers proves
\[
 \operatorname{Var}_R(u)\leq\frac{L(R)(L(R)-1)}4
                      \sum_{e\subset R}|du(e)|^2.
\]
Inside a cluster $L(R)\leq L_k$, and outside all clusters $L(R)\leq n$.
The displayed coefficient lower bounds therefore give
$c\operatorname{Var}_R(u)\leq\mathcal E_{a,R}(u)$.
Singleton residuals are zero, so no zero lower bound is divided by.
Sum over all field/internal coordinates and all partition cells. Original
internal positive-axis edge occurrences are disjoint; omitted inter-cell
and periodic-wrap occurrences have nonnegative energy. Thus the ACTUAL
normalized means
\[
 (C_{\mathcal P}U)_{R,j}=\sqrt{|R|}\,\operatorname{mean}_R U_j
\]
satisfy
\[
 c\|U\|^2\leq\mathcal E_a(U)+c\|C_{\mathcal P}U\|^2
\]
for every cochain, with output dimension $h|\mathcal P|$,
$h=D(N_c^2-1)$. This uses an all-cochain variance estimate, not the invalid
inference from a kernel-only margin to fixed penalty $c$. Means need not
preserve the divergence constraint. The inequality consequently also
holds on that constraint subspace without changing $C_{\mathcal P}$.

If $p=D/2$, $\mu\geq1$, and the original critical mass bound is
$\sum_Xa^{-p}\leq Mn^D/\mu^p$, then each supplied cube contributes
at least $L_k^D/(\beta L_k^2)^p=\beta^{-p}$ to this mass. Hence
$J\leq\beta^pMn^D/\mu^p$, and \eqref{eq:explicit-final-size} gives
\[
 \dim F\leq hA_D(1+2^p\beta^pM)
                   \max\{1,n^D/(1+\mu)^p\}.
\]
This proves the restricted supplied-cube consequence with an explicit
constant and without a fat-object BSP theorem as a black box. It does
not construct the supplied cubes from the critical moment, remove either
scale bound or the strong complement bound, resolve the unrestricted
critical rectangular question, or establish a covariant/continuum gap.
No new originality or independent-discovery claim is made: the compressed
splitting strategy is attributed classical geometry, here derived and
adapted completely for this special class. The earlier separately frozen reviewed source PDF and
its full reports and pending permission remain unchanged.
